Tento krátký průvodce je určen učitelům a administrátorům vysokých škol, kteří chtějí rychle a prakticky začít využívat nástroje umělé inteligence ve výuce; popisuje scénáře, ukázky z praxe, pedagogické přínosy a kroky pro bezpečné a etické zavedení do kurzů. (General background knowledge)
Stručná orientace v pojmech a typech nástrojů:
\begin{itemize}
  \item Nástroje pro učitele: asistenti pro přípravu obsahu, autogradery, generátory cvičení a multimédií. (General background knowledge)
  \item Nástroje pro studenty: adaptivní moduly, konverzační chatboty a systémy průběžné zpětné vazby. (General background knowledge)
\end{itemize}
Hlavní otázky průvodce:
\begin{itemize}
  \item Jak AI podporuje učení a zpětnou vazbu; jak definovat výukové cíle a metriky úspěchu? (General background knowledge)
  \item Jaké jsou požadavky na kvalitu a ochranu dat, souhlas studentů, inkluze a jak prakticky vyhodnocovat dopad? (General background knowledge)
\end{itemize}
Ilustrační příklady z praxe:
\begin{itemize}
  \item Chatboti jako kurzoví asistenti; autogradery pro hodnocení kódu; adaptivní moduly měnící obtížnost podle výkonu. (General background knowledge)
\end{itemize}
Praktické kapitoly obsahují scénáře nasazení s postupem a integrací do LMS, ověřené příklady, tipy pro piloty a škálování, a doporučení pro měření účinnosti, řízení rizik, ochranu údajů a etiku. (General background knowledge)